\section{Conducción inteligente vs. inteligencia corporizada: juego visual e interacción física}

La sección anterior explicó cómo funcionan los datos sintéticos en la conducción autónoma; esta sección aborda por qué la inteligencia corporizada es más difícil. La conducción autónoma no es un problema sencillo, pero tiene una enorme ventaja: el automóvil como plataforma existe desde hace más de cien años, una gran cantidad de vehículos circula por caminos reales y las flotas L2+ pueden enviar de forma continua datos reales de visión, de sensores y de conducción. En cambio, no hay millones de robots humanoides o cuadrúpedos funcionando de manera natural en hogares, fábricas y hoteles, por lo que es muy difícil poner en marcha primero el ciclo virtuoso de datos reales.

Por eso Xie Chen (谢晨) sostiene que «la inteligencia corporizada debe apoyarse primero en una gran cantidad de datos sintéticos y después calibrarse con una pequeña cantidad de datos reales». La conducción autónoma se parece más a un juego visual y de reglas: los caminos, los carriles, los participantes del tránsito y el propio vehículo están relativamente bien definidos. La inteligencia corporizada tiene que lidiar con tazas, popotes, refrigeradores, teclados, telas, líquidos, fricción, colisiones y distintos cuerpos de robot. La complejidad de la interacción física impide que dependa solo de las imágenes y los videos que ya existen en internet.

\begin{figure}[H]
\centering
\includegraphics[width=0.94\textwidth]{av-vs-embodied-data.png}
\caption{Conducción inteligente vs. inteligencia corporizada: la conducción inteligente se inclina hacia el juego visual; la inteligencia corporizada depende más de la interacción física. Diagrama conceptual propio, elaborado a partir de la conversación de 00:16:17--00:32:41.}
\end{figure}

\begin{knowledgebox}{Lectura de la figura: por qué el ciclo virtuoso de datos reales es más fácil en la conducción autónoma}
A la izquierda, la conducción autónoma cuenta con una gran cantidad de vehículos reales, un cuerpo estable y una estructura vial, por lo que los datos reales son la mayor parte y los datos sintéticos sirven principalmente para ampliar la cola larga y los problemas de seguridad. A la derecha, la inteligencia corporizada carece de una población de robots que funcionen de manera natural, y sus tareas abarcan objetos, espacio, contacto y acciones; por eso los datos sintéticos y la simulación deben asumir desde el principio el papel de «construir el mundo».
\end{knowledgebox}

\subsection{El cuello de botella de la interacción física}

Esta subsección explica con más detalle qué dificultad añade exactamente la «interacción física». Los datos visuales abundan en internet, y las fotos y los videos pueden ayudar a generar calles, vehículos, árboles y pavimento; pero los datos de interacción física no son solo apariencia, sino fuerza, fricción, ejes de giro, colisiones, materiales, peso, deformación y recuperación tras un fallo. Sin datos de interacción de bajo nivel, es muy difícil que la IA genere de la nada un comportamiento físico confiable.

La puerta del refrigerador es un buen ejemplo de la entrevista. Para que un robot aprenda a abrir un refrigerador, el recurso de simulación no puede ser solo la carcasa de un refrigerador. Necesita el eje de giro de la puerta, las bisagras, la fuerza del cierre magnético, la fuerza necesaria para abrir la puerta en distintos ángulos, los cuerpos de colisión de los cajones y de la puerta, la distribución de distintos modelos de refrigerador y la fuerza de contacto entre la mano del robot y la puerta. No basta con que visualmente parezca un refrigerador: lo decisivo es que físicamente se pueda interactuar con él.

\begin{figure}[H]
\centering
\includegraphics[width=0.96\textwidth]{physical-interaction-bottleneck.png}
\caption{El cuello de botella de la interacción física: el agarre, el contacto, la retroalimentación de fuerza y las diferencias de material hacen que los datos sean especialmente caros. Diagrama conceptual propio, elaborado a partir de la conversación de 00:16:17--00:32:41.}
\end{figure}

\begin{knowledgebox}{Lectura de la figura: la unidad de los datos corporizados no es el frame, sino la interaction}
Cada variable de la figura puede provocar que la tarea falle: si el punto de agarre se desvía un poco, el objeto se resbala; un material distinto cambia la fricción; los errores del actuador se acumulan; y un cuerpo de colisión impreciso hace que la política entrenada en simulación falle en el mundo real. Para la inteligencia corporizada, si un dato no incluye la acción, el resultado mecánico y la retroalimentación del fallo, es muy difícil que le enseñe al robot a realizar tareas de verdad.
\end{knowledgebox}

\begin{warningbox}{No hay que confundir los datos de video con datos de robot}
Los videos de internet pueden mostrarle al modelo «cómo lo hacen los humanos», pero por lo general no contienen el cuerpo del robot, el estado de las articulaciones, las fuerzas de contacto, los comandos de acción ni las variables controlables del entorno. Sirven para el preentrenamiento de conocimiento del mundo, pero no pueden sustituir directamente los datos de interacción en lazo cerrado que requiere el entrenamiento de robots.
\end{warningbox}

\subsection{La vía de la teleoperación y el ciclo comercial}

Esta subsección responde a una objeción frecuente: si los datos reales son importantes, ¿no se podría primero ganar dinero con robots teleoperados y después recolectar datos? Xie Chen considera que esto podría funcionar en algunos escenarios, pero que como vía general difícilmente avanzaría más rápido que una vía basada principalmente en datos sintéticos. La razón es que la teleoperación primero tiene que generar un valor por el que el cliente esté dispuesto a pagar; de lo contrario, solo es un muestreo de alto costo. Además, la teleoperación entre estados o entre países enfrenta regulación, costos de personal, sindicatos, renta de espacios y complejidad operativa.

El ciclo virtuoso de datos de la conducción autónoma es sólido porque el automóvil ya se vende por sí mismo y el usuario, al manejar, envía los datos de paso; los robots todavía no han llegado a ese punto. Si un robot teleoperado no genera un ROI suficiente en escenarios como aeropuertos, hoteles, restaurantes o comercio minorista, no puede formar el ciclo positivo de «el uso pagado genera datos, los datos mejoran la capacidad y la capacidad trae más uso». El ponente opina que Tesla quizá tenga los escenarios internos de fabricante automotriz y la capacidad organizacional para recorrer parte de esa vía, pero la mayoría de las empresas de robótica no puede simplemente copiarla.

\begin{knowledgebox}{Términos clave: vía real, teleoperación y vía sintética}
\begin{tabularx}{\textwidth}{p{0.20\textwidth}p{0.35\textwidth}X}
\toprule
Vía & Fuente de datos & Restricción principal \\
\midrule
Vía de flota real & Equipos producidos en serie que funcionan de manera natural y envían datos & Requiere que el producto ya sea útil y que la escala de hardware ya exista. \\
Vía de teleoperación & Una persona controla el robot a distancia y se registran las trayectorias de acción & Presión muy alta de personal, escenarios, regulación y ROI. \\
Vía de datos sintéticos & Generación masiva de tareas e interacciones en escenarios de simulación & Requiere Real2Sim de alta calidad, motores de física y sistemas de evaluación. \\
Vía híbrida & Poca calibración con datos reales y ampliación a gran escala con datos sintéticos & La dirección de ingeniería con más probabilidades de volverse dominante actualmente. \\
\bottomrule
\end{tabularx}
\end{knowledgebox}

\subsection{Resumen de la sección}

El cuello de botella de datos de la conducción autónoma se concentra principalmente en la cola larga y en los escenarios de seguridad; el de la inteligencia corporizada llega hasta la propia interacción física. Los robots no cuentan con un ciclo virtuoso de datos reales a gran escala ya establecido, por lo que los datos sintéticos y la simulación física no son un parche, sino el núcleo de la vía inicial.
